% !TEX encoding = UTF-8 Unicode
% !TeX program = xelatex
% !TEX TS-program = XeLaTeX

\documentclass[12pt, a5paper]{article}
\usepackage{mathtools}
\usepackage{unicode-math}
\usepackage{comment}
\usepackage{tikz}
\usepackage[pass]{geometry}
%\geometry{legalpaper, landscape, margin=1in}
\usepackage{fullpage}
\usepackage{setspace}
\usepackage[localise]{xepersian}
\setstretch{1.3}
\settextfont[Scale=1, BoldFont=XB NiloofarBd.ttf]{XB Niloofar.ttf}
\setlatintextfont[Scale=1]{XB Niloofar}%{Junicode} %Junicode
\setdigitfont[BoldFont=XB NiloofarBd.ttf]{XB Niloofar.ttf}%{XB Zar}

\setlength{\parindent}{0ex}

\begin{document}
	
	\begin{tikzpicture}[remember picture,overlay]
	\node[anchor=north,yshift=-1.5cm] at (current page.north) {\rl{به نام خدا}};
	\end{tikzpicture}%
	
	\vspace{-1cm}
	{\centering
		\textbf{کوئیز دوم درس ساختمان داده‌ها و الگوریتم‌ها}\\
		مدرس: مهدی صفرنژاد       - ۴ آبان ۹۸\\
		\raggedleft\small
		نام و نام خانوادگی:\\
		شماره دانشجویی:\\
	}
	\vspace{-0.5cm}
	\پرش‌متوسط\پرخط‌ا\\
	تمام سؤال‌ها جواب کوتاه هستند و توضیح یک یا دو خطی می‌خواهند.
	
	\پاراگراف{سؤال ۱. }
	الگوریتمی از 
	$O(n+k \log(n))$ 
	بیابید که 
	$k$
	بزرگترین اعداد یک آرایه را برگرداند.
	(۱۰ نمره)
	\vspace{3.7cm}
	
	
	\پاراگراف{سؤال ۲.}
	الگوریتمی از 
	$
	O(n)
	$
	ارائه کنید که اعداد مندرج 
	در درخت دودویی جستجو را به ترتیب افزایش مقدار خروجی بدهد.
	(۶ نمره)
		\vspace{3cm}
	\پاراگراف{سؤال ۳.}
	دنباله‌ی اعداد 
	2, 1, 4, 5 و 9
	را به ترتیب از \textbf{راست به چپ} وارد یک درخت قرمز-سیاه بکنید و مرحله به مرحله درخت را رسم کنید.
	(۱۰ نمره)
	
		\vspace{3cm}

	
	\begin{flushleft}
		
		شاد باشین :)
	\end{flushleft}
	
\end{document}